\subsection{从长安到洛阳：两座都城之间的行政距离}
武周选择洛阳并不意味着长安从国家生活中消失。两座都城之间的道路、水运、驿传、仓储和留守官署，使皇帝、官员、家族和物资在不同时间被拉向不同中心。690年的改号需要在洛阳完成仪式与人事安排；关中仍是军事、宗室和供给无法忽略的地区。把这种双中心局面描绘为一个朝廷“迁都”的单一决定，会看不见日常行政如何在两地、沿线州县和众多转运者之间移动。

都城考古所见的道路、坊市、宫城和仓储能够帮助我们把本纪里的抽象地点变成有尺度的空间，却不能由已发掘的区域推出全城人口或市场总量。正史中的朝会、改元与政变偏向宫廷视角；墓志、诗文和宗教题记则让少数迁居者、官员和施主露面。它们共同提示一件事：都城不是一位皇帝的背景，而是把税粮、消息、仪式、教育和人事接合起来的密集基础设施。每当继承冲突爆发，门禁和禁军之所以关键，正因为谁控制都城的文书和通道，谁就更可能让外地先收到自己的版本。

开元时期，玄宗能够在两都之间调度，不表示运输从此没有摩擦。关中对江淮粮食的需要、洛阳的政治地位和沿路州县的劳役都使两都关系成为财政问题。城内的繁华文本多来自能写作的官员、诗人和宾客；它们可说明消费、交游和制度感受，不能代替市场中搬运者、手工业者和临时居住者的生活史。将两都并置，目的不是比较谁更“繁荣”，而是看国家怎样借两处城市获得资源，又怎样因距离、季节和信息延迟受限。

\subsection{钱、布帛与并不整齐的财政世界}
户籍与漕运以外，开元财政还经由布帛、钱币、市场税、土地出产和临时征发进入不同层级。制度志将这些项目排成可读的门类，地方交易和实际输送却常在门类之间移动。农户可能以实物、劳役或折纳回应要求；商人和船户则在市场、驿路与仓场之间寻找可行的交换。我们没有一份足以覆盖帝国每个县的账簿，因此不能用一条法令或一个正史数字计算全部家庭的负担。

钱币窖藏、窑址、仓储遗存和契约能显示某些物资曾经生产、积存或交易，不能单独证明它们代表国家收入或普遍财富。尤其在唐代，窖藏的封藏时间、器物的流通时间和发现的偶然性并不相同。更谨慎的叙事应让不同材料回答不同问题：食货志说明朝廷如何命名税役；文书说明一个地方怎样记录人和物；遗存说明某种物流或生产存在过。三者并读，不会产生一幅整齐的开元经济图，却会让我们看见财政为什么总是既依赖标准化，又离不开地方的变通。

宇文融的括户与裴耀卿的转运可以放在这个不整齐的世界里理解。前者试图扩大可见的征收对象，后者试图让已经征得的资源跨越距离；两者都需要地方官、书吏、船户和仓吏完成。制度带来的收益首先属于能把分散资源接起来的中枢，成本则分布在登记、保管、运输和等待过程中。无法精确计量，并不等于这些成本不存在；恰恰是材料只零星留下它们，才应防止“富庶”一词将差异压平。

\subsection{寺院财产、施主与国家礼仪之间}
武周时期的寺院并非只因《大云经》才进入政治。寺院可以保存经卷、接待旅行者、组织法会、获得施主捐赠，也可能拥有需要管理的房舍、田地、劳力和日常开支。国家给予保护、名额或仪式地位时，影响的是这些可见的组织条件；它不能直接观察每一个人的信仰，也不能决定所有地方寺院如何使用其资源。题记中留下姓名的人往往本就拥有捐施和刻石的能力，不能被扩展为“民众”的统计样本。

对朝廷而言，寺院网络既有可利用的一面，也有需要分类和监管的一面。僧侣的跨区域移动、经卷的传抄和地方施主的名望，能够在国家的官僚渠道之外形成信息和资源联系。对地方家族而言，供养与造像可能是纪念祖先、显示身份、寻求庇护或参与共同体的行动，未必等于接受宫廷的全部政治语言。将这些动机区分开来，并非削弱宗教史，反而使它从“国家利用、社会相信”的二分法中出来。

复唐以后，武周时期的宗教政治并未只留下一个正统判断。年号改变会重写官方叙事，寺院中的文本、图像、器物和施主关系却可能以另一种速度延续。僧传的典范化、碑记的自我呈现和后世史家的褒贬都要求我们校对成书时间与目的。能稳妥说明的，是国家礼仪与佛教组织在这一时期确实相互塑形；不能说明的，是哪一种信仰在所有地区、所有阶层中占据同样位置。

\subsection{妇女、家族与在档案中断裂的生活}
武后、韦后、太平公主和武惠妃使女性在宫廷史中异常显眼，正史也常以此制造道德化的因果链。她们的政治行动当然不能被抹去：继承、诏令、禁军与人事确实经过宫廷家族关系。但将整个时期解释为“女性干政”同样会遮蔽男性宗室、宰相、将领、书吏和地方精英所构成的制度条件，也会使绝大多数不在史书中留名的妇女再次消失。

法律与礼书提供婚姻、继承、户籍和服役的规范词汇；墓志、契约和愿文偶尔让具体家庭的迁居、亲属和财产浮出水面。两种材料都不能直接交代普通妇女如何体验一项政策。墓志常由家族在死亡后撰写，强调德行、门第和追赠；契约只保留一次交易或约定。它们的价值是限制我们对家庭作出过快概括，而不是替沉默者编造完整传记。

当括户、徭役、运输或边防把压力传到乡里，家庭会通过分家、依附、迁移、婚姻、出家或非正式交换来调整。材料通常只在这些选择碰到税役、诉讼、宗教组织或官员时留下痕迹。把这一层放进开元政治史，能够避免将财政写成皇帝和宰相之间的数字游戏：任何重新登记、征调或运粮的决定，都在不同家庭中产生不对称的时间和风险。

\subsection{史料如何让开元看起来过于平滑}
“开元盛世”是一个后世极有力量的概括。它抓住了某些可见的政治稳定、仪式、文化生产与财政能力，却也容易把几十年压成没有摩擦的平面。帝纪按年排列诏令、任免和战报，给人一种中枢连续掌握局势的印象；《通典》《唐会要》将制度分类，使变化看起来像清晰的沿革；诗文和笔记保留城市繁荣的感受。它们都是真实的材料窗口，不能彼此取代。

反过来，不能因为材料有偏向就把开元化为虚构。玄宗即位、括户、封禅、转运、赤岭会盟、登州受袭、南诏册封与天宝改名都有可核的年序和主体。问题在于如何说明它们：封禅不能证明全国富足，册封不能证明稳定统属，户籍奏报不能替代人口普查，边将的军功也不能直接给出控制深度。让这些限度留在叙事里，读者才会看见国家能力并非幻象，却永远是有限、分区并需要维护的能力。

本章从武周到天宝的线索由此闭合。国号的改变、宫廷的政变、寺院的经卷、江淮的粮船、赤岭的盟书和东北、西南的交涉，都是同一政治世界的不同入口。它们相互影响，却不应被压成一种单线因果。下一章的安史危机将在这些已被集中又仍不稳定的资源、道路和人事关系上展开；理解危机之前，必须先理解秩序为何曾经看起来能够维持。
